\chapter{Discussione}

Questa sezione discute alcune funzionalità già descritte precedentemente, con l’obiettivo di illustrare i problemi che persistono nella loro attuale implementazione.
Verranno inoltre esposte alcune possibili soluzioni a questi problemi, unitamente a dei possibili miglioramenti apportabili in futuro.

\section{Problemi}
\subsection{Requisiti per l'utilizzo}
Il problema più evidente nel progetto è che non è una soluzione all-in-one poichè necessita che il sistema abbia Node e Python già installati e visibili nelle variabili di sistema.
Node è necessario a causa del fatto che SvelteKit non permette di compilare in una unica soluzione il sito sviluppato. Il sito sviluppato necessita di un cosiddetto adapter che fungerà da server. 
Python è invece necessario poichè le funzionalità che interagiscono con la libreria win32 utilizzata è implementata in python e tramite Electron viene invocato uno script python per interagire con tale API. 
Oltre a Python però, è necessaria l'installazione di tale api dato che non viene resa disponibile con la sola installazione dell'interprete python. 
Questo pone ovviamente un carico di lavoro eccessivo all'utente dato si necessita che il sistema sia ben configurato affinchè l'applicazione possa funzionare.

\subsection{User Experience}
Come già accennato, l'utente è il responsabile della configurazione corretta del software, processo che potrebbe essere automatizzato. 
Data la complessità nell'automatizzare il processo di mappatura dei programmi installati si è scelto l'approccio drag-and-drop per la categorizzazione dei programmi. 
Questo fornisce una maggiore flessibilità nella scelta, ma va in contrasto con la ragione per cui si è sviluppato tale applicativo, ovvero quella di offrire una interazione più semplice con il dispositivo. 

\subsection{Ambiguità nell'invocazione dei programmi}
A causa della grande mole di software diversi è stato necessario selezionare alcune delle categorie più comuni tra i programmi più popolari. 
Non è possibile chiedere all'assistente di avviare un programma specifico non presente nel dizionario configurato su Actions on Google. 
Un esempio di questo problema è una frase del genere:

\textit{Start IntelliJ IDEA}

L'assistente capisce che deve avviare un programma, ma dato che IntelliJ IDEA non presente nei \textit{types} dell'assistente nè come categoria, nè come alias, non è possibile estrapolare la categoria di software da avviare. 
\'E possibile però assegnare IntelliJ IDEA ad una categoria diversa nel client (ad esempio mappare IntelliJ alla categoria \textit{music}) oppure creare una routine personalizzata in python. 
Entrambe le soluzioni non sono molto intuitive in quanto la prima prevede di avviare un programma utilizzando un altro nome, e la seconda prevede di invocare un programma assegnando un numero al nome dello script che lo invoca.
\'E stato scelto di lasciare però tali possibilità agli utenti che ne potrebbero beneficiare.

\subsection{Deploy}
Il maggiore dei problemi è purtroppo l'effettiva pubblicazione dell'agente al pubblico. 
Attualmente l'agente è disponibile solamente attraverso l'account dello sviluppatore, quindi è possibile provare l'effettivo funzionamento dell'applicazione accedendo con quel profilo. 
Tutte le funzionalità sarebbero comunque disponibili a tutti gli utenti iscritti una volta pubblicato l'agente, ma al momento non è stato possibile rendere disponibile l'agente a tutti. 
Si è scelto di non pubblicare l'agente al momento poichè bisognerebbe stilare una privacy policy e dei termini di servizio compatibili con quelli richiesti da Google, ed al momento questo è fuori dallo scope del lavoro di tesi.

\section{Miglioramenti}
\subsection{Client nativo}
Date le elevate dimensioni dell'applicazione Electron, un miglioramento possibile è quello di reimplementare l'applicazione desktop utilizzando un linguaggio nativo per la costruzione di interfacce grafiche. 
C++ e la libreria wxWidgets sono degli ottimi candidati per il lavoro, e come ulteriore vantaggio avremmo la possibilità di interagire con le API windows per il controllo dei media senza l'utilizzo dell'interprete python eliminando 2 dipendenze.

\subsection{Integrazione di Google Assitant in Windows}
Esiste un progetto open source disponibile su Github che permette di interagire con Google Assistant direttamente da Windows attraverso una interfaccia grafica non ufficiale per Google Assistant.
Questo progetto permetterebbe di eliminare il requisito di utilizzo del dispositivo mobile per impartire comandi, rendendolo opzionale. 
Data la natura open source del progetto, sarebbe possibile integrare la logica di esecuzione dei comandi direttamente in un solo programma.